%! Simplifying Radicals; Adding \& Subtracting — Unit 5, Lesson 5.2 — Warm-Up (Answer Key)
\documentclass[10pt]{article}
\usepackage{algebra2-article}
\usepackage{algebra2-key}   % pulls in -boxes; do NOT also load -boxes
\newcommand{\TallMath}[1]{$\displaystyle #1\rule[-1.4em]{0pt}{3.2em}$}

\begin{document}

\pageheader{Unit 5, Lesson 5.2}{Warm-Up --- Answer Key}
\namedateperiod

\vspace{0.10in}

\begin{notesbox}{Getting Ready: Skills We'll Use Today}
\small These three skills power today's lesson on simplifying and combining radicals. Work quickly.

\vspace{4pt}
\textbf{1. Perfect powers, forwards and backwards.} Evaluate each.
\begin{center}
\renewcommand{\arraystretch}{1.5}
\begin{tabular}{@{}l l l l@{}}
(a) \; $\sqrt{36}=$ \ans{$6$} & (b) \; $\sqrt{81}=$ \ans{$9$} &
(c) \; $\sqrt[3]{27}=$ \ans{$3$} & (d) \; $\sqrt[3]{8}=$ \ans{$2$} \\
\end{tabular}
\end{center}
Now find the \textbf{largest perfect square} that divides each number, and write the number as
that square times what is left over. The first is done for you.
\begin{center}
\renewcommand{\arraystretch}{1.5}
\begin{tabular}{@{}l l l@{}}
$72 = \underline{\ 36\ }\cdot\underline{\ 2\ }$ &
$50 = $ \ans{$25\cdot 2$} &
$48 = $ \ans{$16\cdot 3$} \\
\end{tabular}
\end{center}

\vspace{4pt}
\textbf{2. Combining like terms.} Simplify where you can. If a pair cannot be combined, write
``cannot combine.''
\begin{center}
\renewcommand{\arraystretch}{1.5}
\begin{tabular}{@{}l l l@{}}
(a) \; $3x+7x=$ \ans{$10x$} & (b) \; $3x+7y=$ \ans{cannot combine} &
(c) \; $9m-4m+2p=$ \ans{$5m+2p$} \\
\end{tabular}
\end{center}
What exactly had to match in (a) that did not match in (b)? \ansline{The variable part had to be
identical. In (a) both terms carry $x$, so the coefficients add; in (b) one carries $x$ and the
other $y$, so there is nothing to add --- $3x+7y$ is already as simple as it gets.}

\vspace{4pt}
\textbf{3. Yesterday's exponent rules, one more time.}
\begin{itemize}[leftmargin=1.5em, itemsep=3pt]
  \item Evaluate both: $(4\cdot 9)^{1/2}=$ \ans{$6$} \quad and \quad
        $4^{1/2}\cdot 9^{1/2}=$ \ans{$6$}. \; Equal? \ans{yes}
  \item Write $\dfrac{7}{2}$ as a mixed number: \ans{$3\frac{1}{2}$} \qquad
        Simplify $x^{3}\cdot x^{1/2}=$ \ans{$x^{7/2}$}
\end{itemize}

\end{notesbox}

\begin{teachernote}
Four minutes, then debrief --- every item is a piece of today's lesson wearing a disguise.
\textbf{Item 1} is the whole simplifying algorithm in miniature: today students will read
$72=36\cdot 2$ and pull the $36$ out from under the radical as a $6$. Push on \emph{largest}: a
student who writes $72=4\cdot 18$ is not wrong, but they will finish with $2\sqrt{18}$ and think
they are done. \textbf{Item 2} is the entire adding-and-subtracting rule --- $3\sqrt5+7\sqrt5$
behaves exactly like $3x+7x$, and $3\sqrt5+7\sqrt2$ behaves exactly like $3x+7y$. Get the words
``the variable part had to match'' said out loud now, so that ``the radical part has to match'' is
a two-word edit later. \textbf{Item 3} is the derivation of both properties students will use all
period: the first bullet is $\sqrt{ab}=\sqrt a\sqrt b$ checked numerically before it is stated, and
the second is the surprise ending of Section 4 --- $\sqrt{x^7}=x^{7/2}=x^{3}\cdot x^{1/2}
=x^3\sqrt x$, so simplifying a radical is nothing but rewriting an improper fraction as a mixed
number.
\end{teachernote}

\end{document}
